% Section 8.2, from lectures/L08/README.md: the objectives, the evaluation questions, and the next
% lecture.

\needspace{10\baselineskip}
\section{Review}
\label{c8:sec:review}

\subsection*{What you should be able to do now}
After this chapter, you should:
\begin{itemize}
  \item Have created a concrete subclass \code{Conv} that inherits
    \code{ml::conv_layer::Interface}.
  \item Have implemented \code{feedforward()}, \code{backpropagate()}, and \code{optimize()} for
    \code{Conv}.
\end{itemize}

\needspace{6\baselineskip}
\subsection*{Questions to test yourself}
You should be able to explain:
\begin{enumerate}
  \item What's the difference between \code{ml::conv_layer::Interface} and your concrete class
    \code{Conv}, and why is this split beneficial?
  \item Can you connect the mathematical formulas from the hand-training exercise in
    \secref{c6:sec:exercises} to each computation step in your \code{feedforward()} and
    \code{backpropagate()}?
  \item How does your \code{ml::conv_layer::Conv} compare to the \code{ml::ConvLayer} struct you
    wrote in Chapters~\ref{c6:ch:cnn} and~\ref{c7:ch:layers}: same math, but what's different about
    the design?
  \item Why do \code{Conv} and \code{MaxPool} (Chapter~\ref{c9:ch:maxpool}) share a single
    interface, \code{conv_layer::Interface}?
\end{enumerate}

\needspace{6\baselineskip}
\subsection*{Looking ahead}
Chapter~\ref{c9:ch:maxpool} implements the max pooling layer in the same CNN pipeline.
